\documentclass{article}
\usepackage[margin=25mm]{geometry}
\usepackage{amsmath,amssymb}
\usepackage{xurl}
\usepackage[hidelinks]{hyperref}
% Shared commands for notes in docs/paper-gaps/.

\DeclareUnicodeCharacter{2080}{\ifmmode{}_{0}\else\textsubscript{0}\fi}
\DeclareUnicodeCharacter{2081}{\ifmmode{}_{1}\else\textsubscript{1}\fi}
\DeclareUnicodeCharacter{2082}{\ifmmode{}_{2}\else\textsubscript{2}\fi}
\DeclareUnicodeCharacter{2099}{\ifmmode{}_{n}\else\textsubscript{n}\fi}

\newcommand{\Lean}{\textsc{Lean}}
\ExplSyntaxOn
\NewDocumentCommand{\leanid}{m}{
  \ifmmode
    \textsf{\detokenize{#1}}
  \else
    \group_begin:
    \urlstyle{sf}
    \tl_set:Nn \l_tmpa_tl {#1}
    \tl_replace_all:Nnn \l_tmpa_tl {\_} {_}
    \exp_args:NV \path \l_tmpa_tl
    \group_end:
  \fi
}
\ExplSyntaxOff
\providecommand{\tr}{\operatorname{tr}}
\newcommand{\tnleanIssueBase}{https://github.com/LionSR/TNLean/issues/}
\newcommand{\tnleanPrBase}{https://github.com/LionSR/TNLean/pull/}
\newcommand{\tnleanDocBase}{https://sirui-lu.com/TNLean/docs/}
\newcommand{\ghissue}[1]{\href{\tnleanIssueBase#1}{GitHub issue \##1}}
\newcommand{\ghpr}[1]{\href{\tnleanPrBase#1}{PR \##1}}
\newcommand{\doclink}[1]{\href{\tnleanDocBase#1.html}{\path{#1}}}

% Dirac notation and the identity operator, as in blueprint/src/macros/common.tex.
% \providecommand keeps note-local definitions authoritative.
\usepackage{dsfont}
\providecommand{\ket}[1]{|#1\rangle}
\providecommand{\bra}[1]{\langle #1|}
\providecommand{\braket}[2]{\langle #1 | #2 \rangle}
\providecommand{\ketbra}[2]{|#1\rangle\!\langle #2|}
\providecommand{\Id}{\mathds{1}}
\providecommand{\one}{\mathds{1}}
\providecommand{\C}{\mathbb{C}}
\providecommand{\MN}[1]{M_{#1}(\C)}
\providecommand{\MD}{M_D(\C)}

% Verdict marker: \gapnote{<kind>}{<status>}. Kind is the policy
% classification (clarification, local-correction, scope-restriction,
% unfaithful, false-source, open-gap); status is open, wip (elimination
% underway), resolved, or historical. Machine-read by the site index;
% prints a verdict line.
\newcommand{\gapnote}[2]{\par\noindent\textbf{Verdict:} #1 (#2).\par}

\title{The Missing Diagonal Unit in the Cyclic Wielandt Example}
\author{}
\date{2026-10-29}
\begin{document}
\maketitle
\gapnote{local-correction}{resolved}

\section{The printed step}

The cyclic example in \cite{Sanz2010Quantum}, source
\path{Papers/0909.5347/main.tex:691--700}, has
$S=\sum_{i=0}^{D-1}|i+1\rangle\langle i|$ and
$R=|1\rangle\langle D-1|$, with indices taken modulo $D\ge2$.
Its full-Kraus-rank index is $D(D-1)$.
The upper-bound proof says that at length $N=D(D-1)$ the constraint
\begin{align}
    1+k+l+n(D-1)+mD=N
    \label{eq:spwc10_cyclic_missing_diagonal_padding}
\end{align}
allows arbitrary $0\le k,l<D$ by choosing $n,m\ge0$.
This sentence misses $k=l=D-1$. For example, at $D=3$ that pair
would require $5+2n+3m=6$, which has no nonnegative solutions.
The associated matrix unit is $S^{D-1}RS^{D-1}=|0\rangle\langle0|$.
The example's claimed index and dimension range remain correct.

\section{The correction}

Write $s=k+l$. For $s\le D-2$, choose $n=s+1$ and $m=D-s-2$.
For $D-1\le s\le2D-3$, choose $n=s+1-D$ and $m=2D-s-3$.
Both choices are nonnegative and satisfy~(\ref{eq:spwc10_cyclic_missing_diagonal_padding}).
Hence every matrix unit except $|0\rangle\langle0|$ belongs to $S_N$.
Since $S^N=I$ belongs to the same exact-length span,
\begin{align}
    |0\rangle\langle0|
    =I-\sum_{i=1}^{D-1}|i\rangle\langle i|\in S_N.
    \label{eq:spwc10_cyclic_missing_diagonal_corner}
\end{align}
Equation~(\ref{eq:spwc10_cyclic_missing_diagonal_corner}) supplies the
omitted sentence and proves $S_N=\mathcal M_D$.
It also covers $D=2$.

\section{The lower obstruction}

Every nonzero word containing $R$ reduces to $S^kRS^l$ with length
$1+k+l+n(D-1)+mD$, where $0\le k,l<D$.
At either position $(D-1,0)$ or $(0,1)$, a nonzero entry requires
$k+l=2D-3$. At length $N-1$, the length constraint would read
$1+(2D-3)+n(D-1)+mD=D(D-1)-1$, that is,
\begin{align}
    D(D-1)
    &=(n+m+2)(D-1)+m+1.
    \label{eq:spwc10_cyclic_missing_diagonal_pred_length}
\end{align}
Since $m+1\ge1$, (\ref{eq:spwc10_cyclic_missing_diagonal_pred_length})
gives $(n+m+2)(D-1)<D(D-1)$, and hence
\begin{align}
    n+m+2
    &\le D-1,
    \label{eq:spwc10_cyclic_missing_diagonal_coefficient}\\
    m+1
    &\le n+m+1
     \le D-2.
    \label{eq:spwc10_cyclic_missing_diagonal_remainder}
\end{align}
Substituting (\ref{eq:spwc10_cyclic_missing_diagonal_coefficient}) and
(\ref{eq:spwc10_cyclic_missing_diagonal_remainder}) into the right-hand side
of (\ref{eq:spwc10_cyclic_missing_diagonal_pred_length}) gives
\begin{align}
    (n+m+2)(D-1)+m+1
    &\le(D-1)^2+D-2
     =D(D-1)-1,
    \label{eq:spwc10_cyclic_missing_diagonal_contradiction}
\end{align}
which contradicts (\ref{eq:spwc10_cyclic_missing_diagonal_pred_length}).
The pure shift $S^{N-1}$ has equal entries at the two positions.
Thus every matrix in $S_{N-1}$ has equal entries there, whereas
$|D-1\rangle\langle0|$ does not. This proves $S_{N-1}\ne\mathcal M_D$.
Invertibility of $S$ makes full exact-word span persist to every
larger length; therefore the first full length is $N$.

The formal declarations are in
\path{TNLean/MPS/Examples/CyclicWielandt/}.

\bibliographystyle{alpha}
\bibliography{references}
\end{document}
